%%%%%%%%%%%%%%%%%%%%%%%%%%%%%%%%%%%%%%%%%%%%%%%%%%%%%%%%%%%%%%%%%%%%%%
% Lớp 10 | Bất phương trình bậc 2 một ẩn | Phương trình quy về phương trình bậc hai
% Dạng 0D7?3-1: Phương trình căn
% Số câu: 27  (sinh tự động bởi tach_cau.py; từ nay thầy sửa trực tiếp file này)
%%%%%%%%%%%%%%%%%%%%%%%%%%%%%%%%%%%%%%%%%%%%%%%%%%%%%%%%%%%%%%%%%%%%%%

%%%%%%%%%%%%%%%% TRẮC NGHIỆM 4 PHƯƠNG ÁN %%%%%%%%%%%%%%%%

% nguồn: MaHoaNH_TracNghiem_Lop10
\begin{ex}%[0D7N3-1]%[ma:b0b76481]
	Bình phương cả hai  vế của phương trình $\sqrt{x^2+x+2} = \sqrt{3x+1}$ rồi biến đổi, thu gọn ta được phương trình nào sau đây?
	 \choice
		 {$x^2+x+1=0$}
		{\True $x^2-2x+1=0$} 
		{$x^2-2x-1=0$} 
		{$-x^2+2x+1=0$} 
	 \loigiai{
		Bình phương hai vế phương trình ta  được $x^2+x+2=3x+1 \Leftrightarrow x^2-2x+1=0$.
	 }
\end{ex}

% nguồn: MaHoaNH_TracNghiem_Lop10
\begin{ex}%[0D7N3-1]%[ma:93d068b9]
	Tập nghiệm  của phương trình $\sqrt{2x^2-3x-1}=\sqrt{2 x-4}$ là
	\choice 
		{\True $S=\varnothing$} 
		 {$S=\{1\}$}
		 {$S=\mathbb{R}$}
		{$S=\left\{1;  \dfrac{1}{2}\right\}$}
	\loigiai{Ta có $\sqrt{2x^2-3x-1}=\sqrt{2x-4} \Rightarrow \heva{&2x^2-3x-1=2x-4  \\ &2x-4 \geq 0}\Rightarrow\heva{&2x^2-5x+3=0 \\ &x \geq 2}$\\
		$\Rightarrow\heva{&x=1 \text{ hoặc } x=\dfrac{3}{2} \\ 
		&x \geq 2.}$\\ 
	Vậy phương trình vô nghiệm.} 
\end{ex}

% nguồn: MaHoaNH_TracNghiem_Lop10
\begin{ex}%[0D7N3-1]%[ma:b73a4271]
	Phương trình tổng quát của đường thẳng đi qua  $M(-3; 1)$ và có vectơ pháp tuyến $\overrightarrow{n}=(2; 3)$ là
	\choice 
		 {\True $2x+3y+3=0$}
		 {$2x+3y+5=0$}
		{$3x+2y-9=0$} 
		{$-3x+y+2=0$} 
	\loigiai{ 
		Phương trình của đường thẳng đi qua  $M(-3; 1)$ và có vectơ pháp tuyến $\overrightarrow{n}=(2; 3)$ là\\
		\centerline{$2(x+3)+3(y-1)=0  \Leftrightarrow 2x+3y+3=0$.}
	} 
\end{ex}

% nguồn: MaHoaNH_TracNghiem_Lop10
\begin{ex}%[0D7N3-1]%[ma:cdb16bd0]
	Nghiệm của phương trình $\sqrt{2x-6}=\sqrt{x-2}$ là 
	\choice 
		 {$x=2$}
		{\True $x=4$} 
		{$x=3$} 
		{$x=1$} 
	\loigiai{ 
		Bình phương hai  vế của phương trình đã cho ta được $2x-6=x-2 \Rightarrow x=4$.\\
		Thay $x=4$ vào phương trình đã cho ta thấy $x=4$ thỏa mãn.\\ 
		Vậy  $x=4$ là nghiệm duy nhất của phương trình.
	 }
\end{ex}

% nguồn: MaHoaNH_TracNghiem_Lop10
\begin{ex}%[0D7N3-1]%[ma:732c3722]
	Tính tổng các nghiệm của phương trình $\sqrt{x^2+x+11} =  \sqrt{-2x^2-13x+16}$.
	\choice 
		 {$\dfrac{16}{3}$}
		{$\dfrac{14}{3}$} 
		{\True $-\dfrac{14}{3}$} 
		 {$-\dfrac{16}{3}$}
	\loigiai{ 
		Ta có  $\sqrt{x^2+x+11}=\sqrt{-2x^2-13x+16}$ (1).\\
		Bình  phương hai vế của phương trình (1) ta được $x^2+x+11=-2x^2-13x+16$ (2).\\
		Ta  có $(2) \Leftrightarrow 3x^2+14x-5=0 \Leftrightarrow \hoac{&x=-5\\&x=\dfrac{1}{3}.}$\\
		Thay lần lượt $x=-5$ và $x=\dfrac{1}{3}$ vào phương trình (1) ta thấy  $x=-5$ và $x=\dfrac{1}{3}$ đều thỏa mãn.\\
		Vậy phương trình (1)  có hai nghiệm là $x=-5$ và $x=\dfrac{1}{3}$.\\
		 Nên tổng các nghiệm của phương trình (1) là $-5+\dfrac{1}{3}=-\dfrac{14}{3}$.
	} 
\end{ex}

% nguồn: MaHoaNH_TracNghiem_Lop10
\begin{ex}%[0D7H3-1]%[ma:bac0f135]
	Bình phương cả hai vế của phương trình $\sqrt{x+2}=\sqrt{3x+1}$ rồi biến  đổi, thu gọn ta được phương trình nào sau đây?
	 \choice
		 {$3x-1=0$}
		{$2x+1=0$} 
		{\True  $2x-1=0$}
		{$2x+3=0$} 
	 \loigiai{
		Ta có $\sqrt{x+2}=\sqrt{3x+1}\Leftrightarrow x+2=3x+1\Leftrightarrow  2x-1=0$.
	 }
\end{ex}

% nguồn: MaHoaNH_TracNghiem_Lop10
\begin{ex}%[0D7H3-1]%[ma:4db9e1d9]
	Nghiệm của phương trình  $\sqrt{x^2 + 3x - 2} = \sqrt{x + 1}$ là
	\choice 
		{\True $x = 1$} 
		{$x= -3$} 
		{$\hoac{& x  = -3 \\ & x = 1}$}
		{$ x  \in \emptyset $}
	 \loigiai{
		Điều kiện: $\heva{&x^2 + 3x  - 2 \geq 0\\&x + 1 \geq 0}$.\\
		Ta có $\sqrt{x^2 + 3x - 2} = \sqrt{x + 1} \Leftrightarrow x^2 + 3x - 2 = x + 1 \Leftrightarrow x^2 + 2x - 3 = 0  \Leftrightarrow \hoac{&x = 1\\&x = -3 \;\;\text{(loại)}}$.\\
		Vậy phương trình  đã cho có nghiệm là $x = 1$.
	 }
\end{ex}

% nguồn: MaHoaNH_TracNghiem_Lop10
\begin{ex}%[0D7H3-1]%[ma:2db005ae]
	Tập nghiệm của phương trình $2\sqrt{x^2 - x - 1} = \sqrt{x^2  + 2x + 5}$ là
	 \choice
		{$S = \{1\}$} 
		{$S = \{3\}$} 
		{\True  $S=\{1;3\}$}
		{$S  = \emptyset$}
	 \loigiai{
		$2\sqrt{x^2  - x - 1} = \sqrt{x^2 + 2x + 5}.\quad (1)$\\
		Bình phương hai vế của phương trình $(1)$ ta được phương trình  $4(x^2 - x - 1) = x^2 + 2x + 5$\\
		Ta có $4(x^2 - x - 1) = x^2 + 2x + 5  \Leftrightarrow 3x^2-6x-9=0 \Leftrightarrow \hoac{& x=-1\\& x=3.}$\\
		Thay hai giá  trị $x=-1$ và $x=3$ vào phương trình $(1)$ ta thấy cả hai giá trị đều thõa mãn. Vậy phương trình có hai nghiệm là $x=-1$ và $x=3$.
	 }
\end{ex}

% nguồn: MaHoaNH_TracNghiem_Lop10
\begin{ex}%[0D7H3-1]%[ma:7efbd230]
	Tổng của tất cả các  nghiệm của phương trình $\sqrt{x^2-5x+3}=\sqrt{2x-3}$ là
	 \choice
		 {\True $6$}
		 {$7$}
		 {$1$}
		{$9$} 
	 \loigiai{
	Bình phương hai  vế của phương trình, ta được
	\[x^2-5x+3=2x-3\Leftrightarrow  x^2-7x+6=0\Leftrightarrow \hoac{&x=1\\&x=6.} \]
	Thử lại, ta thấy chỉ có $x=6$ thỏa  mãn.\\
	Vậy tổng các nghiệm  của phương trình đã cho là $6$.
	} 
\end{ex}

% nguồn: MaHoaNH_TracNghiem_Lop10
\begin{ex}%[0D7H3-1]%[ma:522968be]
	Số nghiệm của phương  trình $\sqrt{2x^2+3x-8}=\sqrt{x^2-4}$ là
	 \choice
		 {$2$}
		{\True $1$} 
		 {$3$}
		 {$0$}
	\loigiai{ 
		Ta  có
		\begin{eqnarray*}
			\sqrt{2x^2+3x-8}=\sqrt{x^2-4} &\Leftrightarrow& \heva{&x^2-4 \geq 0 \\  &2x^2+3x-8=x^2-4} \\
			&\Leftrightarrow& \heva{& \hoac{&x \geq 2 \\ &x \leq  -2} \\& x^2+3x-4=0}
			\Leftrightarrow \heva{& \hoac{&x \geq 2 \\ &x \leq -2}  \\ &\hoac{&x=1 \ (\text{loại}) \\ &x=-4 \ (\text{nhận}).}}
		\end{eqnarray*}
	Vậy phương trình đã cho có  $1$ nghiệm.}
\end{ex}

% nguồn: MaHoaNH_TracNghiem_Lop10
\begin{ex}%[0D7H3-1]%[ma:5a576269]
	 Số nghiệm của phương trình $\sqrt{x^2 + x - 42} = \sqrt{2x - 30}$ là
	 \choice
		{\True  $0$}
		 {$1$}
		{$2$} 
		 {vô số}
	\loigiai{ 
		Bình phương  hai vế của phương trình ta được
		\[ 
		x^2 + x - 42=2x  - 30 \Leftrightarrow x^2-x-12=0 \Leftrightarrow
		 \hoac{&x=4\\&x=-3.}
		 \]
		Thay hai giá trị $x=-3$ và $x=4$ vào  phương trình ban đầu ta thấy không có giá trị nào thõa mãn phương trình.\\
		Vậy phương  trình vô nghiệm.
	 }
\end{ex}

% nguồn: MaHoaNH_TracNghiem_Lop10
\begin{ex}%[0D7H3-1]%[ma:d594b7c2]
	Tổng tất cả  các nghiệm của phương trình $\sqrt{x^2 + 2x - 3} = \sqrt{15-5x}$ là
	 \choice
		 {$7$}
		{\True  $-7$}
		{$6$} 
		 {$4$}
	\loigiai{ 
		Bình phương hai vế của phương trình ta  được
		\[ 
		x^2 + 2x -3 = 15-5x \Rightarrow x^2 + 7x-18 = 0\Rightarrow  \hoac{&x=2\\ &x=-9.}
		 \]
		Thử lại ta thấy $x=2$ và $x=-9$ đều thoả  mãn nên tổng các nghiệm của phương trình là
		\[ 
		2+(-9)=  -7.
		 \]
	 }
\end{ex}

% nguồn: MaHoaNH_TracNghiem_Lop10
\begin{ex}%[0D7H3-1]%[ma:e7f45dab]
	Biết phương trình $\sqrt{2x^{2}+3x+1}=2x$ có một nghiệm $x=\dfrac{a+\sqrt{b}}{c}$, với $a$, $b$, $c\in \mathbb{N}^{*}$ và $\dfrac{a}{c}$, $\dfrac{b}{c}$ là các phân số tối giản. Tính  $S=ac-b$.\\
	 \choice
		 {$29$}
		 {$-9$}
		{$5$} 
		 {\True $-5$}
	\loigiai{ 
		\begin{eqnarray*}
			 &&\sqrt{2x^{2}+3x+1}=2x\\
			&\Leftrightarrow&\heva{&2x\geq  0\\&2x^{2}+3x+1=4x^{2}}\\
			 &\Leftrightarrow&\heva{&x\geq 0\\&2x^{2}-3x-1=0}\\
			&\Leftrightarrow&\heva{&x\geq 0\\&\hoac{&x=\dfrac{3-\sqrt{17}}{4}\quad(\text  {loại})\\&x=\dfrac{3+\sqrt{17}}{4}\quad(\text {nhận})}}\\
			&\Leftrightarrow& x=\dfrac{3+\sqrt{17}}{4}. 
		\end{eqnarray*}
		Đối chiếu đề bài ta có  $a=3$; $b=17$; $c=4$.\\
		Vậy $S=ac-b=3\cdot 4-17=-5$. 
	} 
\end{ex}

% nguồn: MaHoaNH_TracNghiem_Lop10
\begin{ex}%[0D7H3-1]%[ma:df0c33ff]
	Nghiệm của phương trình $\sqrt{x^2 - 2x +  4} = \sqrt{2 - x}$ là
	 \choice
		{$x =  \dfrac{1}{2} $}
		{$x = \dfrac{7}{4} $} 
		{$\hoac{& x = \frac{1}{2} \\& x =  \dfrac{7}{4}}$}
		{\True $x \in \emptyset  $}
	\loigiai{ 
		Điều  kiện: $\heva{&x^2 - 2x + 4 \geq 0\\&2 - x \geq 0} \Leftrightarrow \heva{&x \in \mathbb{R}\\&x \leq 2} \Leftrightarrow x \leq 2$.\\
		 Ta có
		 \allowdisplaybreaks
		\begin{eqnarray*}
			\sqrt{x^2 - 2x + 4} = \sqrt{2 - x} 
			&\Leftrightarrow&x^2 - 2x  + 4 = 2 - x\\
			&\Leftrightarrow&x^2 - x + 2 =  0 \;\;\text{(vô nghiệm)}
		\end{eqnarray*}
		Vậy phương trình đã  cho vô nghiệm.
	 }
\end{ex}

% nguồn: MaHoaNH_TracNghiem_Lop10
\begin{ex}%[0D7H3-1]%[ma:0cf98a5b]
	Nghiệm của phương trình $\sqrt{x^2 - 5x + 4} = \sqrt{-2x^2 - 3x +  12}$ là
	 \choice
		{\True $ x =  2$}
		{$ x  = -\dfrac{4}{3} $}
		{$\hoac{& x = 2\\& x  = -\dfrac{4}{3}}$}
		{$ x  \in \emptyset$}
	 \loigiai{
		Điều kiện: $\heva{&x^2 - 5x + 4 \geq 0\\&-2x^2 - 3x + 12 \geq 0}$.\\ 
		Ta có 
		 \allowdisplaybreaks
		\begin{eqnarray*}
			 \sqrt{x^2 - 5x + 4} = \sqrt{-2x^2 - 3x + 12}
			&\Leftrightarrow&x^2 - 5x + 4 =  -2x^2 - 3x + 12\\
			&\Leftrightarrow&3x^2 - 2x - 8  = 0\\
			&\Leftrightarrow&\hoac{&x = 2 \,\text{ (loại)}\\&x = -\dfrac{4}{3}\, \text{ (thoả mãn)}.} 
		\end{eqnarray*}
		Vậy phương trình đã cho có nghiệm là $x  = 2$.
	} 
\end{ex}

% nguồn: MaHoaNH_TracNghiem_Lop10
\begin{ex}%[0D7H3-1]%[ma:4fe298ac]
	Số nghiệm của  phương trình $\sqrt{3x^2-6x+1}=\sqrt{x^2-3}$ là
	\choice 
		 {\True $1$}
		{$0$} 
		 {$3$}
		{$2$} 
	\loigiai{ 
		Bình phương hai vế phương  trình ta được
		 \allowdisplaybreaks
		\begin{eqnarray*}
			 &&3x^2-6x+1=x^2-3
			\\ 
			 &\Leftrightarrow&2x^2-6x+4=0
			\\ 
			 &\Leftrightarrow&\hoac{&x=1\\&x=2.}
		\end{eqnarray*}
		Thử lại chỉ có  $x=2$ thỏa phương trình.\\
		Vậy phương trình  có một nghiệm $x=2$.
	} 
\end{ex}

% nguồn: MaHoaNH_TracNghiem_Lop10
\begin{ex}%[0D7H3-1]%[ma:08a3c510]
	Nghiệm của phương trình  $\sqrt{x^2 - 4x + 4} - \sqrt{(6 - x)(2x - 1)} = 0$ là
	 \choice
		 {$ x = 5 $}
		 {\True $ x = \dfrac{2}{3} $}
		{$\hoac{& x =  \dfrac{2}{3}\\ &x = 5}$}
		{  $ x\in \emptyset$}
	\loigiai{ 
		Điều kiện: $\heva{&x^2 - 4x + 4 \geq 0\\&(6 - x)(2x - 1) \geq 0} \Leftrightarrow \heva{&x \in \mathbb{R}\\&x \in \left(-\infty; \dfrac{1}{2}\right] \cup \left[6; +\infty\right)}  \Leftrightarrow x \in \left(-\infty; \dfrac{1}{2}\right] \cup \left[6; +\infty\right)$.\\
		Ta có 
		\allowdisplaybreaks 
		\begin{eqnarray*}
			\sqrt{x^2 -  4x + 4} - \sqrt{(6 - x)(2x - 1)} = 0
			&\Leftrightarrow&\sqrt{x^2 - 4x + 4}  = \sqrt{(6 - x)(2x - 1)}\\
			&\Leftrightarrow&x^2 - 4x + 4 = (6 - x)(2x -  1)\\
			&\Leftrightarrow&x^2 - 4x + 4 = - 2x^2 + 13x - 6\\ 
			 &\Leftrightarrow&3x^2 - 17x + 10 = 0\\
			&\Leftrightarrow& \hoac{&x  = 5 \;\;\text{(loại)}\\&x = \dfrac{2}{3}.}
		\end{eqnarray*}
		Vậy phương trình đã cho  có nghiệm là $x = \dfrac{2}{3}$.
	} 
\end{ex}


%%%%%%%%%%%%%%%% ĐÚNG/SAI %%%%%%%%%%%%%%%%

% nguồn: MaHoaNH_DungSai_Lop10
\begin{ex}%[0D7N3-1]%[ma:3f9fa9f7]
	Cho phương trình $\sqrt{4x+3}=\sqrt{5x-8}$. 
	\choiceTF 
		 {\True Phương trình có một nghiệm}
		 {\True Phương trình có nghiệm là $x=11$}
		{Phương trình vô  nghiệm}
		{Phương trình có  hai nghiệm phân biệt}
	\loigiai{ 
		Bình phương hai vế của phương trình $\sqrt{4x+3}=\sqrt{5x-8}$ ta  được
		\[4x+3=5x-8\Leftrightarrow x=11\] 
		 Thay giá trị $x=11$ vào phương trình ta thấy thỏa mãn.\\
		Vậy phương trình có  một nghiệm là $x=11$.
	} 
\end{ex}

% nguồn: MaHoaNH_DungSai_Lop10
\begin{ex}%[0D7H3-1]%[ma:a507c802]
	Cho phương trình $\sqrt{x^2-4x-2}=\sqrt{3x^2-8}$. 
	\choiceTF 
		{Phương trình  có hai nghiệm phân biệt}
		{\True Phương trình có một  nghiệm}
		 {Phương trình có nghiệm $x=1$}
		{\True Phương  trình có nghiệm $x=-3$}
	\loigiai{ 
		 Bình phương hai vế của phương trình $\sqrt{x^2-4x-2}=\sqrt{3x^2-8}$ ta được
		\begin{eqnarray*}
			 && x^2-4x-2=3x^2-8\\
			&\Leftrightarrow& -2x^2-4x+6=0\\ 
			&\Leftrightarrow&  \hoac{&x=1\\&x=-3.}
		\end{eqnarray*}
		Thay lần lượt các giá  trị của $x$ vào phương trình ta thấy chỉ có $x=-3$ thỏa mãn.\\
		Vậy phương trình có một nghiệm là $x=-3$. 
	} 
\end{ex}

% nguồn: MaHoaNH_DungSai_Lop10
\begin{ex}%[0D7V3-1]%[ma:c73fbf7c]
	Cho phương  trình $\sqrt{5x^2-8x+2}=\sqrt{x^2+2}$ (${*}$). Khi đó
	\choiceTF 
		{\True $x^2+2\geq 0$  đúng $\forall x \in \mathbb{R}$}
		{Bình phương hai  vế ta được $4x^2-3x=0$}
		{\True Phương trình $\left(^{*}\right)$ có 2 nghiệm} 
		{Tổng các  nghiệm của phương trình $\left(^{*}\right)$ bằng $0$}
	\loigiai{ 
		\begin{itemchoice}
			\itemch $x^2\geq 0$ đúng $\forall x \in \mathbb{R}$  nên $x^2+2\geq 0$ đúng $\forall x \in \mathbb{R}$.
			 \itemch Bình phương hai vế ta được $5x^2-8x+2=x^2+2\Leftrightarrow 4x^2-8x=0$
			\itemch Bình phương hai vế (*) và rút  gọn ta được $\hoac{&x=0\\
			&x=2.}$\\ 
			Vậy tập  nghiệm của phương trình là $S=\{0; 2\}$.
			\itemch  Tổng các nghiệm của (*) là $0+2=2$.
		\end{itemchoice}
	} 
\end{ex}

% nguồn: MaHoaNH_DungSai_Lop10
\begin{ex}%[0D7V3-1]%[ma:90248127]
	Cho phương trình $\sqrt{x^2 - 4x - 5}  = \sqrt{2x^2 + 3x + 1} \quad (1)$.
	\choiceTF 
		{Bình phương hai vế của phương trình $(1)$ ta được phương trình $x^2 - 7x  + 6 = 0$}
		{\True $x=-1$ là nghiệm của phương trình $(1)$} 
		{Tổng các nghiệm của phương trình $(1)$ là $-5$} 
		{\True Phương  trình sao có nhiều hơn $1$ nghiệm}
	\loigiai{ 
		Bình phương hai vế của phương  trình $(1)$ ta được
		 \[
		x^2 - 4x - 5 = 2x^2 + 3x + 1 \Leftrightarrow x^2+7x+6 = 0 \Leftrightarrow  \hoac{&x=-1\\ &x=-6.}
		\] 
		Thử lại ta thấy $x=-1$ và $x=-6$ đều thoả  mãn. Vậy
		\begin{itemchoice}
			\itemch Bình phương hai vế và thu gọn  ta được phương trình $x^2 + 7x + 6 = 0$.
			\itemch Phương trình đã cho  có nghiệm $x=-1$.
			\itemch Tổng các nghiệm  của phương trình $(1)$ là $-1+(-6)=-7$.
			\itemch Phương trình $(1)$ có $2$ nghiệm $-1$ và $-6$. 
		\end{itemchoice}
	 }
\end{ex}

% nguồn: MaHoaNH_DungSai_Lop10
\begin{ex}%[0D7V3-1]%[ma:fb9787c5]
	Cho phương trình $(x+1)\left( \sqrt{x+4} - \sqrt{-x^2 + 4x + 14} \right) = 0 \quad (1)$.  Khi đó
	 \choiceTF
		{Phương trình $(1)$  không có nghiệm bé hơn $4$}
		{\True Phương trình $(1)$  có $3$ nghiệm phân biệt}
		{\True Các  nghiệm của phương trình $(1)$ đều là số nguyên}
		{Tổng bình phương các nghiệm của  phương trình $(1)$ bằng $29$}
	\loigiai{ 
		Ta  có
		\[ 
		(1) \Rightarrow \hoac{&x+1 = 0 \\ & \sqrt{x+4} - \sqrt{-x^2 + 4x  + 14} = 0}
		\Rightarrow \hoac{&x=-1\\ &x+4 =  -x^2 + 4x + 14}
		\Leftrightarrow \hoac{&x=-1\\ &x^2 - 3x -  10 = 0.}
		\] 
		Từ đó ta giải được các nghiệm $x=-1$, $x=-2$,  $x=5$.\\
		Thử lại ta  thấy $x=-1$, $x=-2$ và $x=5$ đều thoả mãn. Vậy
		\begin{itemchoice}
			\itemch Phương trình có nghiệm  $x=-1 < 4$.
			\itemch Phương trình $(1)$ có $3$ nghiệm $x=-1$, $x=-2$  và $x=5$.
			\itemch Ta có $-1$; $-2$; $5$ đều  là các số nguyên.
			\itemch Ta có $(-1)^2 + (-2)^2  + 5^2 = 30.$
		\end{itemchoice}
	 }
\end{ex}


%%%%%%%%%%%%%%%% TRẢ LỜI NGẮN %%%%%%%%%%%%%%%%

% nguồn: MaHoaNH_TraLoiNgan_Lop10
\begin{bt}%[0D7H3-1]%[ma:caa0bc56]
	Tổng các nghiệm của phương trình $\sqrt {x^2+3x-2}=\sqrt {1+x}$ là 
	 \shortans{1}
	\loigiai{Ta  có
		\begin{eqnarray*}
			\sqrt{x^2+3x-2}=\sqrt{1+x}&\Rightarrow  &x^2+3x-2=1+x\\
			&\Rightarrow  &x^2+2x-3=0\\
			&\Rightarrow&\hoac{&x=1\\&x=-3.} 
		\end{eqnarray*}
	Thay $x=1$ và $x=-3$ vào phương trình bài ra ta  thấy $x=1$ thoả mãn. Vậy tổng các nghiệm của phương trình là $1$.}
\end{bt}

% nguồn: MaHoaNH_TraLoiNgan_Lop10
\begin{bt}%[0D7H3-1]%[ma:8d37428d]
	 Tính tổng các nghiệm của phương trình $\sqrt{3x^2-4x-1}=\sqrt{2x^2-4x+3}$.
	 \shortans {0}
	\loigiai{Ta có 
		\begin{eqnarray*}
			&&\sqrt{3x^2-4x-1}=\sqrt{2x^2-4x+3}\\ 
			&\Rightarrow&x^2-4x-1=2x^2-4x+3 \text  { (bình phương hai vế)}\\
			&\Leftrightarrow& x^2=4\\ 
			 &\Leftrightarrow& \hoac {&x=2 \\&x=-2.}
		\end{eqnarray*}
		Thay  lần lượt hai giá trị này của $x$ vào phương trình đã cho, ta thấy cả hai đều thỏa mãn.\\
		Vậy tổng các nghiệm của phương trình đã cho là  $2+(-2)=0$.
	 }
\end{bt}

% nguồn: MaHoaNH_TraLoiNgan_Lop10
\begin{bt}%[0D7H3-1]%[ma:26bf51b4]
	Gọi $S$ là tập nghiệm của phương trình $\sqrt{x^2+3x-2}=\sqrt{1+x}$.  Tổng tất cả các phần tử của $S$ bằng bao nhiêu?
	\shortans{1} 
	 \loigiai{
		Bình phương hai vế của phương trình ta có 
		\allowdisplaybreaks 
		\begin{eqnarray*}
			x^2+3x-2=1+x&\Leftrightarrow& x^2+2x-3=0\\ 
			&\Leftrightarrow&\hoac{&x=1\\&x=-3.} 
		\end{eqnarray*}
		Lần lượt thế $x=1$ và $x=-3$ vào phương trình ban đầu  chỉ có $x=1$ thỏa mãn.\\
		Do đó, phương trình có  tập nghiệm là $S=\{1\}$.
	 }
\end{bt}


%%%%%%%%%%%%%%%% TỰ LUẬN %%%%%%%%%%%%%%%%

% nguồn: MaHoaNH_TuLuan_Lop10
\begin{bt}%[0D7H3-1]%[ma:cba79be7]
	Phương trình $\sqrt{x^2+2x-3}=\sqrt{-2x^2+5}$ có bao nhiêu  nghiệm?
	\loigiai{Ta  có
		\begin{eqnarray*}
			 &&\sqrt{x^2+2x-3}=\sqrt{-2x^2+5}\\
			&\Rightarrow&x^2+2x-3=-2x^2+5 \text  { (bình phương hai vế)}\\
			 &\Leftrightarrow& 3x^2+2x-8=0\\
			 &\Leftrightarrow&\hoac {&x=\dfrac{4}{3} \\ &x=-2.}
		\end{eqnarray*}
		Thay lần lượt hai giá trị này của $x$ vào phương trình đã cho, ta  thấy $x=\dfrac{4}{3}$ thỏa mãn.\\
		Vậy phương trình đã cho  có một nghiệm là $x=\dfrac{4}{3}$.
	 }
\end{bt}

% nguồn: MaHoaNH_TuLuan_Lop10
\begin{bt}%[0D7V3-1]%[ma:67651a93]
	Giải phương trình $\sqrt{x^2 + 2x + 4} = \sqrt{2-x}$. 
	\loigiai{ 
		 Bình phương hai vế của phương trình đã cho ta được
		\begin{eqnarray*}
			x^2 + 2x + 4 = 2-x&\Leftrightarrow&  x^2+3x+2=0\\
			&\Leftrightarrow& x=-1\text{ hoặc  } x=-2.
		\end{eqnarray*}
		 Thay lần lượt các giá trị trên vào phương trình đã cho, ta thấy $x=-1$ và $x=-2$ thỏa mãn.\\
		Vậy nghiệm của phương  tình đã cho là $-1$ và $-2$.
	} 
\end{bt}
